% ==============================================================================
% EPÍGRAFE (opcional; manual, cap. 2, item 8)
% - Não tem título: a palavra "Epígrafe" NÃO aparece na página.
% - Vem após os agradecimentos, em folha própria, sem número impresso.
% - É uma citação direta e segue a NBR 10520:2023: aspas duplas, autor em
%   minúsculas, ano e página (quando houver). Texto traduzido pelo autor
%   leva a indicação "tradução nossa".
% - A obra citada DEVE constar nas Referências, na edição efetivamente
%   consultada. Não usar frases de autoria incerta, comuns na internet.
% - Preferir citações relacionadas ao tema ou ao espírito da pesquisa.
% - A NBR 14724 também admite epígrafes na abertura de cada capítulo.
% ==============================================================================
\begin{epigrafe}
  ``Se vi mais longe, foi por estar sobre os ombros de gigantes.''

  \vspace{0.5em}
  \hfill (Newton, 1676, tradução nossa)
\end{epigrafe}

% Exemplo de referência correspondente (adaptar à edição consultada):
% NEWTON, Isaac. [Carta a Robert Hooke, 5 fev. 1676]. In: TURNBULL, H. W.
% (ed.). The correspondence of Isaac Newton. Cambridge: Cambridge University
% Press, 1959. v. 1, p. 416.